%%
%% This is file `sample-sigconf.tex',
%% generated with the docstrip utility.
%%
%% The original source files were:
%%
%% samples.dtx  (with options: `all,proceedings,bibtex,sigconf')
%% 
%% IMPORTANT NOTICE:
%% 
%% For the copyright see the source file.
%% 
%% Any modified versions of this file must be renamed
%% with new filenames distinct from sample-sigconf.tex.
%% 
%% For distribution of the original source see the terms
%% for copying and modification in the file samples.dtx.
%% 
%% This generated file may be distributed as long as the
%% original source files, as listed above, are part of the
%% same distribution. (The sources need not necessarily be
%% in the same archive or directory.)
%%
%%
%% Commands for TeXCount
%TC:macro \cite [option:text,text]
%TC:macro \citep [option:text,text]
%TC:macro \citet [option:text,text]
%TC:envir table 0 1
%TC:envir table* 0 1
%TC:envir tabular [ignore] word
%TC:envir displaymath 0 word
%TC:envir math 0 word
%TC:envir comment 0 0
%%
%% The first command in your LaTeX source must be the \documentclass
%% command.
%%
%% For submission and review of your manuscript please change the
%% command to \documentclass[manuscript, screen, review]{acmart}.
%%
%% When submitting camera ready or to TAPS, please change the command
%% to \documentclass[sigconf]{acmart} or whichever template is required
%% for your publication.
%%
%%
\documentclass[manuscript, review]{acmart}
%%
%% \BibTeX command to typeset BibTeX logo in the docs
\AtBeginDocument{%
  \providecommand\BibTeX{{%
    Bib\TeX}}}

%% Rights management information.  This information is sent to you
%% when you complete the rights form.  These commands have SAMPLE
%% values in them; it is your responsibility as an author to replace
%% the commands and values with those provided to you when you
%% complete the rights form.
\setcopyright{acmlicensed}
\copyrightyear{2026}
\acmYear{2026}
\acmDOI{ImStillLearningLaTeXDoIKeepThis}

%%
%%  Uncomment \acmBooktitle if the title of the proceedings is different
%%  from ``Proceedings of ...''!
%%
%%\acmBooktitle{Woodstock '18: ACM Symposium on Neural Gaze Detection,
%%  June 03--05, 2018, Woodstock, NY}
\acmISBN{978-1-4503-XXXX-X/2018/06}


%%
%% Submission ID.
%% Use this when submitting an article to a sponsored event. You'll
%% receive a unique submission ID from the organizers
%% of the event, and this ID should be used as the parameter to this command.
%%\acmSubmissionID{123-A56-BU3}

%%
%% For managing citations, it is recommended to use bibliography
%% files in BibTeX format.
%%
%% You can then either use BibTeX with the ACM-Reference-Format style,
%% or BibLaTeX with the acmnumeric or acmauthoryear sytles, that include
%% support for advanced citation of software artefact from the
%% biblatex-software package, also separately available on CTAN.
%%
%% Look at the sample-*-biblatex.tex files for templates showcasing
%% the biblatex styles.
%%

%%
%% The majority of ACM publications use numbered citations and
%% references.  The command \citestyle{authoryear} switches to the
%% "author year" style.
%%
%% If you are preparing content for an event
%% sponsored by ACM SIGGRAPH, you must use the "author year" style of
%% citations and references.
%% Uncommenting
%% the next command will enable that style.
%%\citestyle{acmauthoryear}


%%
%% end of the preamble, start of the body of the document source.
\begin{document}

%%
%% The "title" command has an optional parameter,
%% allowing the author to define a "short title" to be used in page headers.
\title{Screen Scrubbing: Multimodal Interaction in Video Editing}

%%
%% The "author" command and its associated commands are used to define
%% the authors and their affiliations.
%% Of note is the shared affiliation of the first two authors, and the
%% "authornote" and "authornotemark" commands
%% used to denote shared contribution to the research.
\author{Davis Garrett}
\affiliation{%
  \institution{Colorado State University}
  \city{Fort Collins}
  \state{Colorado}
  \country{USA}
}

\author{Riley Bowling}
\affiliation{%
  \institution{Colorado State University}
  \city{Fort Collins}
  \state{Colorado}
  \country{USA}
}

\author{Connor Schwab}
\affiliation{%
  \institution{Colorado State University}
  \city{Fort Collins}
  \state{Colorado}
  \country{USA}
}


%%
%% By default, the full list of authors will be used in the page
%% headers. Often, this list is too long, and will overlap
%% other information printed in the page headers. This command allows
%% the author to define a more concise list
%% of authors' names for this purpose.
\renewcommand{\shortauthors}{Garrett et al.}

%%
%% The abstract is a short summary of the work to be presented in the
%% article.
\begin{abstract}
    Lorem Ipsum || Study has yet to finish
\end{abstract}

%%
%% Keywords. The author(s) should pick words that accurately describe
%% the work being presented. Separate the keywords with commas.
\keywords{Multimodal, Interaction, Virtual Reality, VR, Video, Video Editing}

\received{13 April 2026}
%%\received[revised]{12 March 2009}

%%
%% This command processes the author and affiliation and title
%% information and builds the first part of the formatted document.
\maketitle

\section{Introduction}
A timeline is the centerpiece of any video editing software. It’s the collective figure showing ones footage, sound, effects and the primary point of interaction with such. Despite the importance it holds over the software the methods of interaction with the tool remain in a gully. However, with the growing advancements in the field of virtual reality (VR) it’s expected to see these mediums see crossing users. Beyond the capabilities of traditional flat screen displays virtual reality enables both a more customizable 3D environment, but also changes the means and methods of interaction. In particular, the growing adoption of gestures based interaction in VR applications could support a new means of navigation within video editing.

Gestures provide a more natural means of interaction in comparison to classic mouse and keyboard. Instead of relegating movement and commands to the combination of two plus input devices, gestures enable users to rely on physical hand and finger movement, simplifying controls. Additionally these factors work alongside the added environment the application is situated in, serving both virtual and augmented reality (AR).

Additionally the use of VR provides the opportunity for expanded multimodal interaction without the increased need for external hardware (IE cameras, sensors) beyond what current market headsets already provide. In tandem with gestures interaction methods such as gaze and audio could work to strengthen one another, particularly in ways that continue the naturalistic control scheme gestures encourage. This function is especially ideal in the context of video editing, where the media is most likely a 2D representation being bent to the metaphor of a 3D space, a common problem in VR when dealing with user interaction.

As of writing, the question of video editing, in the context or virtual reality, remains solely focused on the aspects of VR film, or 360 degree videography. Programs designed specifically to handle media tailored for audiences using VR. Even so, all popular software engaging with such does so through the more familiar guise of the 2D editor. In this paper, we looked to improve understandings within how the classic format of video editing, focusing on interactions with the timeline element, could be improved or otherwise substituted in a 3D virtual space using multimodal interactions.


\section{Related Work}

\subsection{360 Video Editing}

In the space of virtual reality video is a complicated subject, most obvious due to the increased range of viewpoints a creator must address. Unlike 2D screens VR enables projection anywhere in relation to a viewer, with full immersion requiring 360 degree. In video editing, this grows complex, with a high degree of stitching needed to sell the illusion. This is the process of taking multiple cameras or video feeds and systematically adjusting them to expand the visible content, generally taking anywhere from 2 to upwards of 12+ feeds \cite{LeirpollJarle2017VVE}. The process is itself relatively new in the consumer market, and software that supports it is lax.

Araújo et al. \cite{deCastroAraujoGabriel2025A3VP} in their 2025 works were unable to find any open source program pertaining to 360 degree video editing, with the study being their own endeavor in creating a practical application for the purpose of filling this void, and documenting the findings along the way. Their final product, 360EAVP, addresses the issues of the timeline among their core issues, specifically the alignment and navigation of content. Due to the stitching required, the UI becomes increasing cluttered as more footage was incrementally added, a problem seen commonly in non-liner editors \cite{GriffinRuairi202166DI}.

While the bulk of the study was focused on the aspects of streamlining hardware issues, such as when users have headsets with differing FOV, the techniques within static editing remain useful within the context of the timeline. The technique used for addressing struggles when managing multiple videos across the same period in particular shows potential in how it could translate to menus in 3D spaces, with pointing and gaze help to create a fixed location, or better adjust position. In their findings, while not quantiative, users reported a far better experience with 360EAVP, and critically positive workflow,  pivotal when professionals report a mere minute of video takes a minimum of an hour to properly edit \cite{RickyTelg2021VE}. Given the prolonged nature editing takes, maintaining a user’s comfort is a requirement second to few. While VR can improve certain fatigue, extended use can begin to toll a user physically \cite{10918867}.

\subsection{Multimodal Interaction}

Attempting to move beyond basic mouse and keyboard means the reduction of input devices, but rather than a pure loss of control it serves as a potential for greater expansion. In 2013 Zhao et al. \cite{ZhaoShengdong2013SIMM} studied the impacts of mixed modalities across stationary and mobile devices, focusing heavily on visual and audio. Focusing on their first experiment environment, the collaboration was interested in feedback as to how extra stimuli would impact a subject performance in simple menu navigation. The conclusions of their findings indicated both an improvement in response time when multimodal cues were used, and likewise user feedback veered towards a preference in the additional visual / audio.

Many multimodal projects follow similarly, with how the combination of interaction means increase user response, both qualitatively and quantitatively. One of the earliest examples of such, the Put-That-There study \cite{cdi_proquest_ebookcentralchapters_5927840_6_21} continues to support, through it’s own findings, that well balanced multimodality leads to a general increase in a user’s performance. This extended across audio, 2D and 3D environments with users reporting that it allowed for quicker, simpler recovery from mistakes, give a more personal way to interact.

\subsection{2D to 3D VR Interaction}

Core to our study is the ability to bend 2D interfaces to the metaphor of a 3D space. Unlike prior works, this subject has seen fierce scrutiny. Starting as early as 2009 research was being done on how the projection of 2D systems could be bent to 3D works in the public market. Schöning et al. \cite{SchoningJohannes2009BIwI} were among early adopters seeking to approach how interaction of flat elements could be best super imposed in VR spaces. The study approaches the issue of degrees of freedom (DoF), and how most 2D applications only support two as it works best for such medium. Thus, the translation into 3D environments, where the DoF often ranges from three to six, is a game of trade offs as the increased DoF does not, by function, increase usability, even hindering applications not intended for such.

In this light the focus becomes how to best translate 2D UI(s) to functioning, 3D elements. This breaks down into two primary components; the presentation of content, and means of interacting with content. Taking 2D elements into the 3D visually is seen as the simpler of the two. A 2024 study \cite{LuzsaRobert2026Tpei} researching VR and MR (mixed reality) content, among other issues, delved into how well 2D content transferred directly into 3D digital spaces, and found that even with utmost minimal changes the presentation of content remained almost identical. Users found little difference in how they experienced flat media regardless if it was on a monitor or headset, but that was only half the story. Virtual reality turned to show a trade off in how users experienced media, not the content itself but the quality in which is was presented. The virtual space enabled a series of changes that were otherwise impossible in a 2D setting, such as the removal of distractions and giving greater freedom of movement to view, which were positively remarked. However, due to the hardware limitations of headsets, especially regarding pixel density, content could appear disjointed, blurry, or otherwise changed in a manner unintended from the original presentation. Ultimately the study found users generally preferred positive polarity content in VR and MR over direct translation, and equally to the 2D original.

More complex is how users interact with this new content in VR. Song et al. released their study in 2022 using the foil of keyboard macros to understand how well basic and familiar concepts could feasibly translate \cite{9995649}. The 2D interface they worked with was that of a Qwerty keyboard layout used on mobile devices, and seeing if; one, users could reasonably adapt to the element in 3D, and two, if a new function for special characters could be added that also sees quick adoption. Despite the familiar presentation of the keyboard, only 11 of the 16 participants found the interaction familiar, and barring that a mere seven of the group found it natural to use. Frustrations grew around many technical aspects of the project, such as misread inputs, but some noted a general annoyance with VR altogether, a known issue as more immersive applications tend to elevate higher emotional response \cite{10443512}.

Technical struggles of using 3D inputs for 2D content are consistent throughout virtually all studies discussing the subject. Even among the most basic of controls, such as clicking a button or hovering a desired selection, there are often problems overlooked as they never occur in the 2D space. Woojoo Kim and Shuping Xiong in 2024 \cite{KimWoojoo2024TMTo} extensively studied these forms of basic interaction, focused on dominant techniques such as raycasting. They found that while certain methods fair better than others, all suffered from issues revolving around precision, as the interaction was meant for 2D mouse, leading to difficulty using smaller elements and problems with the shaky nature of the human body.

\section{Methodology}

\subsection{Equipment}

The primary device was an Oculus Quest 2 serving as the headset, and likewise means of interacting with the program in VR. This was connected directly to a Predator Helios 300 serving as the PC to both facilitate the program, collect and calculate data, and run 2D elements of the experiment. The laptop was running a Windows 10 environment, with 16 GB of ram and GTX 1660Ti GPU. The open board CPU was a 4 core i7 9th Gen. The VR program itself was made using Unity 6.3. The 2D program was designed with a modified version of the open source program kdenlive to properly track and log data from participants. All data was captured directly by the program and store later for further processing.

\subsection{Participants}

Participants were selected at random, consisting of 16 total participants with an even gender split of eight males and eight females. All participants were students studying at Colorado State University, the location the study took place. The age range varied between 20 and 2X. The ultimate goal with our participants was to have an even spread of users who were both familiar and unfamiliar with video editing, as to asses the potential learning curve and possible influences time invested has on impact. Participants were informed the study would last roughly ten to fifteen minutes and involve tasks both sitting on a computer and sitting in a VR / virtual 3D space, with little physical strain or movement to be expected.

\subsection{Procedure}

When a participant would enter the room they would be informed of the prior before being handed three forms. The first of the three reiterated the experiments procedure, the types of data which would be collected, and possible side effects they may have to VR. Signing was to confirm both their consent and right to leave anytime they feel unwell, uncertain, or any other concerning detail. The second form consisted of a questionnaire inquiring about their basic information (age, gender, major), as well as any prior experience they may have. This was focused on two subject: do you have have history with video editing (or related subjects, if so what)? Do you have any history with virtual reality (or related subject, if so what)? A third form was to be retain for after the experiment, which asked about the participants experience with the programs, which they preferred, and what they felt was more or less difficult to accomplish in VR.

Once the first two forms were finished and collected, they would then be assigned a group. For the first group the participant was seated at the computer and given a brief demo of the program, kdenlive including basic controls and layout, and then given three sequential tasks. The first task was to find and navigate to five colors in order: red (4), orange (2), yellow (1), green (3), blue (5). Time started once the timeline moved and ended when the center bar touched the final color. The second task was to remove sections with the color red. This used a second video in which four sections required deletion. Time started once the timeline moved and ended when the final deletion was made. The final task was to rearrange a video such that the colors matched the following: purple, orange, red, white, blue. Time started once the timeline moved and ended with the final edit.

When all computer tasks were finished, the participant would then be given an Quest 2 and asked to wear it. The participant would then be given a likewise tutorial of the VR video editor, shown the basic controls and layout. This demo was extended slightly compared to the 2D portion to help acclimate participants to how the headset read their gestures. Once ready, the interaction would begin, again asking for the same tasks as before, however using different videos requiring a new solution to the previous. The first task, in order: red (5), orange (2), yellow (3), green (1), blue (4), and the third task requiring: orange, white, blue, red, purple. All experiments continued to start once the timeline moved and ending on their respective condition.

For the remaining group, the order of operations was reversed. Those in group two were instead started on the VR editor before moving to the 2D environment, with all prior tasks and instructions remaining the same. Each of the eight in their respective group were also broken into three sub groups. One would experience the VR editor without any additional interaction, the second had auditory cues for commands and inputs, and the third had notifications for inputs and additional information.

Once all tasks for both mediums were complete the participant would be asked to then fill the third form to the best of their abilities. Data from both the 2D and 3D programs were collected and compiled for future use. Once finished, the participant was offered food and drink for their time before they left.

\subsection{Data Collection}

All programs were designed to start and stop recording when certain conditions were met. For all tasks the start condition was when the timeline element of the respective video editor was moved, indicating the participant has begun to engage with the task. Each of the three tasks then checks for their respective conditions.

The first video noted the exact time a participant touched a new color in the sequence, requiring at least five frames before accepting, with the timer ending once the final color in the sequence was reached, then dumping all the information into a log.

The second video checked when certain chunks of the video were still present. When a piece was deleted the time was logged, though due to limitations the exact one removed was not. The timer stopped and log dumped when all required blocks were no longer present.

The final video suffered a similar issue to the second where the exact position of the videos blocks were unable to be tracked. Due to this the actual program stopped tracking once the final frame of the video matched the desired color, with researchers closely observing to make sure no participant accidentally trigger this oversight (IE deleting other colors, misordering cuts, etc). No participant suffered this issue, thus all data collected was used.

All respective data was then condensed together, matching groups to their respective members and double checking for any egregious outliers that were missed to avoid possible issues stemming from an issue during testing.

%%
%% The next two lines define the bibliography style to be used, and
%% the bibliography file.
\bibliographystyle{ACM-Reference-Format}
\bibliography{sample-base}


%%
%% If your work has an appendix, this is the place to put it.
\appendix


\end{document}
\endinput
%%
%% End of file `sample-sigconf.tex'.
